\documentclass[11pt,letterpaper]{article}
\usepackage[T1]{fontenc}
\usepackage{amsmath,amsthm,amssymb}
\usepackage[margin=1.05in]{geometry}
\usepackage{microtype}
\setlength{\emergencystretch}{1em}
\usepackage{xcolor}
\definecolor{paperlink}{RGB}{0,0,114}
\usepackage{hyperref}
\hypersetup{colorlinks=true,linktoc=all,
  linkcolor=paperlink,citecolor=paperlink,urlcolor=paperlink,pdfstartview=Fit,
  pdftitle={A five-generated family of numerical semigroups with first Betti number 2w+2},
  pdfauthor={Ben Keene}}
\usepackage{enumitem}
\setlist[itemize]{leftmargin=2em,itemsep=0.25em,topsep=0.35em}
\setlist[enumerate]{leftmargin=2.25em,itemsep=0.25em,topsep=0.35em}
\setlength{\parindent}{1.25em}
\setlength{\parskip}{0pt}
\usepackage{fancyhdr}
\pagestyle{fancy}
\fancyhf{}
\fancyhead[L]{\small A five-generated family with first Betti number $2w+2$}
\fancyhead[R]{\small\thepage}
\setlength{\headheight}{14pt}
\numberwithin{equation}{section}
\setcounter{tocdepth}{1}
\date{}

% Draft markers
\newcommand{\red}[1]{\textcolor{red}{#1}}

% Theorem environments
\theoremstyle{plain}
\newtheorem{theorem}{Theorem}[section]
\newtheorem{proposition}[theorem]{Proposition}
\newtheorem{lemma}[theorem]{Lemma}
\newtheorem{corollary}[theorem]{Corollary}
\newtheorem{conjecture}[theorem]{Conjecture}
\theoremstyle{remark}
\newtheorem{remark}[theorem]{Remark}

% Notation
\newcommand{\NN}{\mathbb{N}}
\newcommand{\mm}{\mathfrak{m}}
\newcommand{\Ap}{\operatorname{Ap}}
\newcommand{\NF}{\mathcal{A}}
\newcommand{\len}[1]{\lvert #1\rvert}

\title{A five-generated family of numerical semigroups\\ with first Betti number $2w+2$}
\author{Ben Keene\thanks{\red{[affiliation]}}}

\begin{document}

\maketitle
\thispagestyle{plain}

\begin{abstract}
For integers $k\ge 3$ and $q\ge k-1$ and $m=2kq+1$, we show that the toric ideal of the numerical semigroup $\langle m,m+1,m+k,m+2k-1,m+2k\rangle$ over any field is minimally generated by $4k+2$ binomials. Thus the first Betti number of these semigroups is $\beta_1=2w+2$, where $w=2k$ is the width, and within this family the linear bound $\beta_1\le 2w$ valid for four-generated numerical semigroups is exceeded by exactly two. The proof writes down an explicit presentation and reduces every monomial to a power of the multiplicity variable times one of $m$ normal forms, which are indexed by the residues modulo~$m$. Nonredundancy follows from projection to a single weight together with a monomial specialisation that separates the two kinds of relations, and, among the relations homogeneous in the ordinary grading, from projection to a single ordinary degree. A graded argument through $I/\mm I$ then gives the lower bound over every generating set. The theorem has been formally verified in Lean~4 with Mathlib.
\end{abstract}

\section{Introduction}\label{sec:intro}

Let $S=\langle n_0,\dots,n_{e-1}\rangle\subseteq\NN$ be a numerical semigroup minimally generated by $n_0<\dots<n_{e-1}$, so that $e$ is its embedding dimension, $n_0$ its multiplicity and $w=n_{e-1}-n_0$ its width. Over a field $K$, the semigroup ring $K[S]$ is the quotient of $K[x_0,\dots,x_{e-1}]$ by the toric ideal $I_S$, the kernel of the map $x_i\mapsto t^{n_i}$. The first Betti number $\beta_1(S)$ is the minimal number of generators $\mu(I_S)$ of this ideal, the least cardinality of a set generating it, that is, the number of relations in a minimal presentation of~$S$. Already in embedding dimension four it is unbounded, a result of Bresinsky recalled in~\cite{CMS24}; the width is the invariant expected to control it.

\paragraph{Width bounds.}
Herzog and Stamate conjectured that the defining ideal of the tangent cone of $K[S]$ needs at most $\binom{w+1}{2}$ generators~\cite[Conjecture~2.1]{HS14}. As noted in~\cite{CMS24}, this implies $\beta_1(S)\le\binom{w+1}{2}$, the form stated by Caviglia, Moscariello and Sammartano~\cite[Conjecture~1.2]{CMS24}. They extended it to all Betti numbers~\cite[Conjecture~1.3]{CMS24}, proved a first bound depending on the width alone, and settled both conjectures for four-generated semigroups of width at least $40$~\cite[Theorem~1.5]{CMS24}. In embedding dimension four the quadratic bound is far from the truth. The author has shown that $\beta_1(S)\le 2w$ for every minimally four-generated numerical semigroup of width~$w$; this bound is formally verified in Lean, in the registry entry~\cite{Palomar4gen}. It is sharp for every $w\ge 4$, by an informal companion result of that entry which is not formalized. The question addressed here is what replaces $2w$ in embedding dimension five.

\paragraph{The result.}
We exhibit a two-parameter family of five-generated semigroups that exceeds the four-generator bound by exactly two.

\begin{theorem}\label{thm:F}
Let $K$ be a field, let $k\ge 3$ and $q\ge k-1$ be integers, let $m=2kq+1$, and let $S=\langle m,m+1,m+k,m+2k-1,m+2k\rangle$. Then the toric ideal $I_S\subseteq K[a,b,c,d,e]$ satisfies $\mu(I_S)=4k+2$.
\end{theorem}

Since $S$ is minimally five-generated of width $w=2k$ (Section~\ref{sec:family}), the theorem says $\beta_1(S)=2w+2$. The proof gives an explicit generating set of $4k+2$ binomials, of which $2k+2$ are homogeneous in the ordinary grading and $2k$ are not. We are not aware of a prior statement of the theorem. We propose the following bound for embedding dimension five.

\begin{conjecture}\label{conj:W}
Every minimally five-generated numerical semigroup $S$ of width $4\le w\le 40$ satisfies $\beta_1(S)\le 2w+2$.
\end{conjecture}

The conjecture is supported by an exhaustive computation for every width $4\le w\le 20$ over all multiplicities $2\le m\le w^2+w$, which covers every multiplicity by the periodicity theorem of Conaway et al.~\cite[Theorem~4.9]{CGHOPWW18}. The upper limit $40$ is a choice, not a mathematical boundary: it is the range for which a certified exhaustive computation over the periodicity window is planned, and it matches the width threshold of the four-generator theorem~\cite[Theorem~1.5]{CMS24}. By Theorem~\ref{thm:F}, if the conjecture holds, its bound is attained at every even width $6\le w\le 40$.

\paragraph{Why a new presentation is needed.}
The offsets $0,1,k,2k-1,2k$ of our generators from the multiplicity contain no four-term arithmetic progression when $k\ge 3$. So the family is neither generated by an arithmetic sequence nor by an almost arithmetic one (all generators but one in arithmetic progression), and presentation formulas for those classes do not apply. Shifting all generators by the width preserves the size of a minimal presentation once the multiplicity exceeds~$w^2$~\cite[Theorem~4.9]{CGHOPWW18}; this reduces each set of offsets to finitely many multiplicities, but it gives no formula, and our family has $m\equiv 1\pmod w$ with multiplicities $m\ge w^2/2-w+1$ both below and above~$w^2$.

\paragraph{Strategy.}
Grade $R=K[a,b,c,d,e]$ by the weights $m,m+1,m+k,m+2k-1,m+2k$, so that $I_S$ is spanned by differences of monomials of equal weight. The generating set consists of $2k+2$ relations between monomials of equal length and $2k$ relations $N_r$, one in each weight $(q+2)m+r$ with $0\le r<2k$, between monomials of lengths $q+1$ and $q+2$. Generation is proved by a reduction to normal forms. There are exactly $m$ normal forms, monomials free of $a$ whose weights represent the $m$ residues modulo $m$ once each, and every monomial is congruent modulo the relations to a power of $a$ times a normal form; induction on the weight of the part free of~$a$ makes this reduction terminate. Hence two monomials of equal weight are congruent, and the relations generate.

Nonredundancy is proved by projection. Projecting a putative redundancy to the weight of $N_r$ leaves only multiples of the relations of equal length. The monomial specialisation $\psi$ into $K[s,t]$, in which the exponent of $s$ records the length of a monomial and the exponent of $t$ its excess weight over $m$ times its length, kills those relations but not~$N_r$, because $\psi(N_r)=s^{q+1}t^r(t^m-s)$. The relations of equal length are separated from one another by projection to a single length, after which, modulo $(a,c^2,bd)$, those of length $k$ have disjoint supports. Nonredundancy of one generating set does not by itself bound the size of the others; the lower bound over every generating set comes from the dimension of $I/\mm I$ (Lemma~\ref{lem:graded}).

\section{The family and its presentation}\label{sec:family}

Throughout the paper $k\ge 3$ and $q\ge k-1$ are integers, $m=2kq+1$, and
\begin{equation}\label{eq:S}
S=\langle m,\ m+1,\ m+k,\ m+2k-1,\ m+2k\rangle\,.
\end{equation}
Since $S$ contains the coprime integers $m$ and $m+1$, it is a numerical semigroup. All five generators lie in the interval $[m,2m)$ because $2k<m$, while a sum of two or more generators is at least~$2m$; hence none of them is a sum of the others, and $S$ is minimally generated by the five displayed integers. Its multiplicity is~$m$ and its width, the difference between the largest and the smallest minimal generator, is $w=2k$.

\paragraph{Weights and offsets.}
Let $K$ be a field, of any characteristic, and let $R=K[a,b,c,d,e]$. We give the variables the weights
\begin{equation*}
\nu(a)=m,\quad \nu(b)=m+1,\quad \nu(c)=m+k,\quad \nu(d)=m+2k-1,\quad \nu(e)=m+2k\,,
\end{equation*}
and extend $\nu$ additively to monomials. The toric ideal of $S$ is the kernel $I=\ker\varphi$ of the $K$-algebra map $\varphi\colon R\to K[t]$ sending each variable $x$ to $t^{\nu(x)}$, so that $\varphi(M)=t^{\nu(M)}$ for every monomial~$M$; Theorem~\ref{thm:F} concerns its minimal number of generators~$\mu(I)$.

For a monomial $M$ we write $\len{M}$ for its ordinary degree, which we call its length, and
\begin{equation*}
\delta(M)=\nu(M)-m\len{M}
\end{equation*}
for its offset, so that $\nu(M)=m\len{M}+\delta(M)$: the weight of a monomial is determined by its length and its offset, and this is how weights are compared below. Offsets are additive, and the variables $a,b,c,d,e$ have offsets $0,1,k,2k-1,2k$. A polynomial is homogeneous of weight $n$ if all of its monomials have weight~$n$; this grades $R=\bigoplus_{n\ge 0}R_n$, and $I$ is a homogeneous ideal because $\varphi$ maps $R_n$ into $Kt^n$. We also use the ordinary grading of $R$, and call a polynomial homogeneous of length $n$ if all of its monomials have length~$n$.

\paragraph{The relations.}
We shall see that two kinds of binomials generate~$I$. The first kind consists of the $2k+2$ binomials
\begin{equation}\label{eq:H}
\begin{gathered}
P_1=c^2-bd,\qquad P_2=bd-ae,\qquad P_3=b^k-a^{k-1}c,\qquad P_4=d^k-ce^{k-1},\\
F_i=b^{k-1-i}ce^{i}-a^{k-1-i}d^{i+1},\qquad G_i=b^{k-1-i}e^{i+1}-a^{k-2-i}cd^{i+1}\,,
\end{gathered}
\end{equation}
where $0\le i\le k-2$. We call this set $H$. The second kind consists of the $2k$ binomials $N_r=U_r-V_r$, $0\le r\le 2k-1$, where
\begin{equation}\label{eq:N}
\begin{aligned}
&U_r=cd^{k-1-r}e^{q+1+r-k}, && V_r=a^{q+2-r}b^{r} && (0\le r\le k-1),\\
&U_{k+s}=d^{k-1-s}e^{q+2+s-k}, && V_{k+s}=a^{q+1-s}b^{s}c && (0\le s\le k-1)\,.
\end{aligned}
\end{equation}
We call this set $N$. In $N_r$ the monomial $U_r$ has offset $m+r$ and $V_r$ has offset~$r$ (Lemma~\ref{lem:relations}): the relation trades an excess offset of $m$ for one more factor, which is why its two sides have lengths $q+1$ and $q+2$. Every exponent is a nonnegative integer: $q+1+r-k\ge r$ and $q+2+s-k\ge s+1$ because $q\ge k-1$, while $q+2-r\ge q+3-k\ge 2$ and $q+1-s\ge q+2-k\ge 1$. At the boundary $q=k-1$ the monomials $U_r$ and $U_{k+s}$ become $cd^{k-1-r}e^{r}$ and $d^{k-1-s}e^{s+1}$. Let $J\subseteq R$ be the ideal generated by $H\cup N$, and put $D=(q+2)m$, the least weight of an element of~$N$.

\begin{lemma}\label{lem:relations}
Each element of $H\cup N$ is the difference of two distinct monomials of the same weight; in particular $J\subseteq I$. The two monomials of each element of $H$ have the same length, namely $2$ for $P_1,P_2$ and $k$ otherwise, and their common weights are
\begin{equation}\label{eq:Hweights}
\begin{gathered}
\nu(P_1)=\nu(P_2)=2m+2k,\qquad \nu(P_3)=km+k,\qquad \nu(P_4)=km+k(2k-1),\\
\nu(F_i)=km+(2k-1)(i+1),\qquad \nu(G_i)=km+k+(2k-1)(i+1)\,.
\end{gathered}
\end{equation}
For $0\le r\le 2k-1$ the monomial $U_r$ has length $q+1$ and offset $m+r$, the monomial $V_r$ has length $q+2$ and offset $r$, and $\nu(N_r)=D+r$.
\end{lemma}

\begin{proof}
Since $\nu(M)=m\len{M}+\delta(M)$, two monomials have the same weight as soon as they have the same length and the same offset, and we compare these. For $P_1,\dots,P_4$ the lengths agree by inspection and the offsets are $2k=1+(2k-1)$, $2k=0+2k$, $k=0+k$ and $k(2k-1)=k+2k(k-1)$. For $F_i$ both monomials have length $k$, and the offsets $(k-1-i)+k+2ki$ and $(2k-1)(i+1)$ agree. For $G_i$ both have length $k$, and the offsets $(k-1-i)+2k(i+1)$ and $k+(2k-1)(i+1)$ agree. This gives~\eqref{eq:Hweights}.

For $0\le r\le k-1$ the monomial $U_r$ has length $1+(k-1-r)+(q+1+r-k)=q+1$ and offset
\begin{equation*}
k+(2k-1)(k-1-r)+2k(q+1+r-k)=2kq+1+r=m+r\,,
\end{equation*}
so $\nu(U_r)=(q+1)m+m+r=D+r$, while $V_r$ has length $q+2$, offset $r$ and weight $D+r$. For $r=k+s$ with $0\le s\le k-1$, the monomial $U_{k+s}$ has length $(k-1-s)+(q+2+s-k)=q+1$ and offset $(2k-1)(k-1-s)+2k(q+2+s-k)=2kq+1+k+s=m+r$, while $V_{k+s}$ has length $q+2$ and offset $s+k=r$. In each binomial the two monomials are distinct, since their lengths or their supports differ. Finally, $\varphi(M-M')=t^{\nu(M)}-t^{\nu(M')}=0$ whenever $\nu(M)=\nu(M')$, so every element of $H\cup N$ lies in~$I$.
\end{proof}

The weights of $N$ fill the band $D,D+1,\dots,D+2k-1$ of $2k$ consecutive integers, one relation in each; Section~\ref{sec:nonred} shows that the elements of $H$ can meet this band only at the boundary $q=k-1$.

\section{Normal forms}\label{sec:normal}

Section~\ref{sec:generation} uses the relations of Section~\ref{sec:family} as rewriting rules on monomials free of~$a$. The rewritings stop exactly at the following monomials, which are divisible by none of the monomials rewritten there. Let $\NF$ be the set of monomials of the four shapes
\begin{equation}\label{eq:NF}
\begin{aligned}
&\text{(A1)}\ \ b^{i}e^{j} \quad (i\ge 1,\ i+j\le k-1), &\qquad& \text{(A2)}\ \ b^{i}ce^{j} \quad (i\ge 1,\ i+j\le k-2),\\
&\text{(A3)}\ \ d^{t}e^{j} \quad (t\le k-1,\ t+j\le q), && \text{(A4)}\ \ cd^{t}e^{j} \quad (t\le k-1,\ t+j\le q-1),
\end{aligned}
\end{equation}
with all exponents nonnegative integers. Elements of $\NF$ are called normal forms. None of them involves~$a$; the families (A1) and (A2) avoid~$d$, the families (A3) and (A4) avoid~$b$, and $c$ occurs at most once. Their offsets are
\begin{equation}\label{eq:NFoffsets}
\begin{gathered}
\delta(b^ie^j)=i+2kj,\qquad \delta(b^ice^j)=k+i+2kj,\\
\delta(d^te^j)=2k(t+j)-t,\qquad \delta(cd^te^j)=k+2k(t+j)-t\,.
\end{gathered}
\end{equation}
The point of the definition is that these offsets are exactly the residues modulo~$m$, each once.

\begin{proposition}\label{prop:normal}
The offset map $W\mapsto\delta(W)$ is a bijection from $\NF$ onto $\{0,1,\dots,m-1\}$. In particular $\NF$ has exactly $m$ elements, and the weights $\nu(W)=m\len{W}+\delta(W)$ of distinct normal forms lie in distinct residue classes modulo~$m$.
\end{proposition}

\begin{proof}
Write an integer $n\ge 0$ as $n=2kh+\rho$ with $h\ge 0$ and $0\le\rho\le 2k-1$. Since $m-1=2kq$, the integers $0\le n\le m-1$ are those with $\rho=0$ and $0\le h\le q$, or with $\rho\ne 0$ and $0\le h\le q-1$. We determine the pairs $(\rho,h)$ attained by each family in~\eqref{eq:NFoffsets}.
\begin{itemize}
\item (A1) $\delta(b^ie^j)=2kj+i$ with $1\le i\le k-1$, so $\rho=i$ and $h=j$, and the constraint reads $0\le h\le k-1-\rho$.
\item (A2) $\delta(b^ice^j)=2kj+(k+i)$ with $1\le i\le k-2$, so $\rho=k+i$ and $h=j$, and $0\le h\le 2k-2-\rho$.
\item (A3) For $t=0$ we get $\rho=0$ and $h=j$ with $0\le h\le q$. For $1\le t\le k-1$ we have $2k(t+j)-t=2k(t+j-1)+(2k-t)$, so $\rho=2k-t\in[k+1,2k-1]$ and $h=t+j-1$, and the constraints read $2k-1-\rho\le h\le q-1$.
\item (A4) For $t=0$ we get $\rho=k$ and $h=j$ with $0\le h\le q-1$. For $1\le t\le k-1$ we have $k-t+2k(t+j)$, so $\rho=k-t\in[1,k-1]$ and $h=t+j$, and $k-\rho\le h\le q-1$.
\end{itemize}
In every case the exponents of $W$ are recovered from $(\rho,h)$ and the family, so it remains to check that, for each value of $\rho=n\bmod 2k$, the ranges of $h$ contributed by the families are disjoint and their union is the required range. The table lists them; a range whose upper end is below its lower end is empty.
\begin{center}
\renewcommand{\arraystretch}{1.25}
\begin{tabular}{lll}
$\rho=n\bmod 2k$ & normal forms & range of $h$\\
\hline
$\rho=0$ & $e^h$ & $0\le h\le q$\\
$\rho=k$ & $ce^h$ & $0\le h\le q-1$\\
$1\le\rho\le k-1$ & $b^{\rho}e^h$ & $0\le h\le k-1-\rho$\\
 & $cd^{k-\rho}e^{h-(k-\rho)}$ & $k-\rho\le h\le q-1$\\
$k+1\le\rho\le 2k-1$ & $b^{\rho-k}ce^h$ & $0\le h\le 2k-2-\rho$\\
 & $d^{2k-\rho}e^{h-(2k-1-\rho)}$ & $2k-1-\rho\le h\le q-1$
\end{tabular}
\end{center}
For $1\le\rho\le k-1$ the two ranges are consecutive and cover $[0,q-1]$, because $k-1-\rho\le k-2\le q-1$. For $k+1\le\rho\le 2k-1$ they are again consecutive and cover $[0,q-1]$, because $2k-1-\rho\le k-2\le q-1$; at $\rho=2k-1$ the first range is empty. Hence the value $\rho=0$ carries $q+1$ normal forms and every other value of $\rho$ carries~$q$; the offsets are distinct, and they fill $\{0,\dots,m-1\}$, a set of $(q+1)+(2k-1)q=m$ elements. The final assertion holds because $\nu(W)\equiv\delta(W)\pmod m$.
\end{proof}

\section{Generation}\label{sec:generation}

We show that $J=I$ by reducing every monomial modulo~$J$ to a power of~$a$ times a normal form. Rewriting by a relation of Lemma~\ref{lem:relations} preserves the weight, so weight alone cannot measure progress; what does decrease is the weight of the part of the monomial free of~$a$. Most rewritings produce a positive power of~$a$, which we split off; of the two that do not, $be^{k-1}\to cd^{k-1}$ and $d^k\to ce^{k-1}$, at most one is applied before a power of $a$ appears or a normal form is reached.

Write $f\equiv g$ for $f-g\in J$. If $u-v$ is an element of $H\cup N$, or its negative, and $u$ divides a monomial~$M$, then $M\equiv (M/u)\,v$, and the two sides have the same weight. We use also the consequences $c^2\equiv ae$ (from $P_1+P_2=c^2-ae$) and $cd^k\equiv c^2e^{k-1}\equiv ae^k$ (from $cP_4$ and $c^2-ae$).

\begin{lemma}\label{lem:reduction}
For every monomial $M\in R$ there are an integer $\ell\ge 0$ and a normal form $W\in\NF$ such that $M\equiv a^{\ell}W$. In particular $\nu(M)=\ell m+\nu(W)$.
\end{lemma}

\begin{proof}
Writing $M=a^pM'$ with $M'$ free of~$a$, it suffices to treat $M'$, since $M'\equiv a^{\ell}W$ implies $M\equiv a^{p+\ell}W$. We argue by strong induction on the weight of an $a$-free monomial $M=b^{\beta}c^{\gamma}d^{\eta}e^{\varepsilon}$, assuming the statement for all $a$-free monomials of smaller weight. Call a congruence $M\equiv a^pM_1$ with $p\ge 1$ and $M_1$ free of~$a$ a \emph{descent}. A descent finishes the proof: $\nu(M_1)=\nu(M)-pm<\nu(M)$, so $M_1\equiv a^{\ell}W$ by induction, and $M\equiv a^{p+\ell}W$. We show that $M$ is a normal form, admits a descent, or is congruent to an $a$-free monomial of the same weight that is a normal form or admits a descent.

\emph{Step 1: three descents.} If $\gamma\ge 2$, then $c^2\equiv ae$ gives a descent. If $\beta\ge 1$ and $\eta\ge 1$, then $P_2$ rewrites $bd$ as $ae$. If $\beta\ge k$, then $P_3$ rewrites $b^k$ as $a^{k-1}c$, and $k-1\ge 1$. We may therefore assume $\gamma\le 1$, $\beta\le k-1$ and $\beta\eta=0$.

\emph{Step 2: monomials containing~$b$.} Let $\beta\ge 1$, so that $\eta=0$ and $M=b^{\beta}c^{\gamma}e^{\varepsilon}$, and put $i=k-1-\beta$, so that $0\le i\le k-2$. If $\gamma=1$ and $\beta+\varepsilon\le k-2$, then $M$ is a normal form of type~(A2). If $\gamma=1$ and $\beta+\varepsilon\ge k-1$, then $\varepsilon\ge i$, so $b^{\beta}ce^{i}$ divides $M$ and $F_i$ rewrites it as $a^{\beta}d^{i+1}$, a descent. If $\gamma=0$ and $\beta+\varepsilon\le k-1$, then $M$ is of type~(A1). If $\gamma=0$ and $\beta+\varepsilon\ge k$, then $\varepsilon\ge i+1$, so $b^{\beta}e^{i+1}$ divides~$M$ and $G_i$ rewrites it as $a^{\beta-1}cd^{i+1}$. For $\beta\ge 2$ this is a descent. For $\beta=1$ it is the rewriting $be^{k-1}\to cd^{k-1}$, which gives $M\equiv cd^{k-1}e^{\varepsilon-k+1}$, an $a$-free and $b$-free monomial treated in Step~3.

\emph{Step 3: monomials without~$b$.} Let $\beta=0$, so $M=c^{\gamma}d^{\eta}e^{\varepsilon}$ with $\gamma\le 1$. If $\eta\ge k$ and $\gamma=1$, then $cd^k\equiv ae^k$ gives a descent. If $\eta\ge k$ and $\gamma=0$, then $P_4$ gives $M\equiv cd^{\eta-k}e^{\varepsilon+k-1}$, which has $\gamma=1$ and is treated by the present step. It remains to consider $\eta\le k-1$.
\begin{itemize}
\item If $\gamma=0$ and $\eta+\varepsilon\le q$, then $M$ is of type~(A3). If $\gamma=0$ and $\eta+\varepsilon\ge q+1$, put $s=k-1-\eta\in[0,k-1]$. Then $U_{k+s}=d^{\eta}e^{q+1-\eta}$ divides $M$, and $N_{k+s}$ rewrites it as $V_{k+s}=a^{q+1-s}b^{s}c$, a descent because $q+1-s\ge q+2-k\ge 1$.
\item If $\gamma=1$ and $\eta+\varepsilon\le q-1$, then $M$ is of type~(A4). If $\gamma=1$ and $\eta+\varepsilon\ge q$, put $r=k-1-\eta\in[0,k-1]$. Then $U_r=cd^{\eta}e^{q-\eta}$ divides $M$, and $N_r$ rewrites it as $V_r=a^{q+2-r}b^{r}$, a descent because $q+2-r\ge 2$.
\end{itemize}
At the boundary $q=k-1$ and $\eta=k-1$ the divisor in the last case is $U_0=cd^{k-1}$, with no factor~$e$. At most one rewriting without descent occurs: each of the two leads to a monomial with $\gamma=1$, to which neither Step~2 nor the rewriting by $P_4$ in Step~3 applies, so it is a normal form or admits a descent. This completes the induction. The weight identity holds because every congruence used preserves the weight.
\end{proof}

The second ingredient is a description of $I$ that does not refer to the family. It is the standard fact that a toric ideal is spanned by binomials.

\begin{lemma}\label{lem:binomials}
As a $K$-vector space, $I$ is spanned by the binomials $M-M'$ with $M,M'$ monomials of the same weight.
\end{lemma}

\begin{proof}
These binomials lie in $I$ because $\varphi(M-M')=t^{\nu(M)}-t^{\nu(M')}$. Conversely, let $f\in I$ and write $f=\sum_{n}f_n$ with $f_n$ homogeneous of weight~$n$. Since $\varphi(f_n)$ equals the sum of the coefficients of $f_n$ times~$t^n$, and these powers of $t$ are distinct, $\varphi(f)=0$ forces every coefficient sum to vanish. Fix $n$ with $f_n=\sum_M\lambda_MM\ne 0$ and a monomial $M_n$ of weight~$n$. Then
\begin{equation*}
f_n=\sum_M\lambda_M(M-M_n)+\Bigl(\sum_M\lambda_M\Bigr)M_n=\sum_M\lambda_M(M-M_n)\,,
\end{equation*}
a combination of binomials of the required kind.
\end{proof}

\begin{proposition}\label{prop:generation}
The binomials $H\cup N$ generate the toric ideal: $J=I$.
\end{proposition}

\begin{proof}
By Lemma~\ref{lem:relations}, $J\subseteq I$. For the converse, by Lemma~\ref{lem:binomials} it suffices to show $M-M'\in J$ whenever $\nu(M)=\nu(M')=n$. Lemma~\ref{lem:reduction} gives $M\equiv a^{\ell}W$ and $M'\equiv a^{\ell'}W'$ with $n=\ell m+\nu(W)=\ell'm+\nu(W')$. Then $\nu(W)\equiv\nu(W')\pmod m$, so $W=W'$ by Proposition~\ref{prop:normal}, hence $\ell=\ell'$, and $M\equiv M'$.
\end{proof}

\begin{corollary}\label{cor:apery}
The Ap\'ery set $\Ap(S;m)=\{s\in S: s-m\notin S\}$ equals $\{\nu(W): W\in\NF\}$, and these $m$ weights lie in distinct residue classes modulo~$m$.
\end{corollary}

\begin{proof}
Every element of $S$ is the weight of some monomial, hence by Lemma~\ref{lem:reduction} has the form $\ell m+\nu(W)$ with $\ell\ge 0$ and $W\in\NF$. If $\ell\ge 1$, then subtracting $m$ stays in~$S$, so every element of $\Ap(S;m)$ is a weight $\nu(W)$. Conversely, suppose $\nu(W)-m\in S$ for some $W\in\NF$. Then $\nu(W)-m=\ell'm+\nu(W')$ with $W'\in\NF$, so $W'=W$ by Proposition~\ref{prop:normal}, and $-m=\ell'm$, which is impossible. The last assertion is Proposition~\ref{prop:normal}.
\end{proof}

\section{Nonredundancy}\label{sec:nonred}

A generating set is nonredundant if none of its elements lies in the ideal generated by the others. We prove that $H\cup N$ is nonredundant. Two gradings separate the relations. The weight $\nu$ isolates each $N_r$: the relations of $N$ occupy $2k$ consecutive weights, while a nonconstant multiplier has weight at least~$m$. The relations of $H$, which are homogeneous of length $2$ or~$k$, are separated from each other by projection to a single length together with reduction modulo $(a,c^2,bd)$. What links the two kinds is the monomial map
\begin{equation*}
\psi\colon R\to K[s,t],\qquad a\mapsto s,\quad b\mapsto st,\quad c\mapsto st^{k},\quad d\mapsto st^{2k-1},\quad e\mapsto st^{2k}\,,
\end{equation*}
which records length and offset, $\psi(M)=s^{\len{M}}t^{\delta(M)}$.

\subsection{Weights and the projection~$\psi$}

When a putative redundancy is projected to the weight of an element $f$, the coefficient of $N_j$ becomes homogeneous of weight $\nu(f)-D-j$. To see that it is then a scalar, we need to know which elements of $H$ have weight below $D+m$.

\begin{lemma}\label{lem:Hweights}
Every element of $H$ has weight less than $D+m$. If $q\ge k$, every element of $H$ has weight less than~$D$. If $q=k-1$, exactly two elements of $H$ have weight in $[D,D+m)$, namely $\nu(G_{k-2})=D$ and $\nu(P_4)=D+k-1$.
\end{lemma}

\begin{proof}
By~\eqref{eq:Hweights}, a relation $f$ of length $k$ has weight $km+\delta_f$, where $\delta_f\le k(2k-1)$ is the common offset of its two monomials, the largest offset being that of~$P_4$; indeed $\delta(G_i)=3k-1+(2k-1)i\le 2k^2-2k+1$ and $\delta(F_i)\le(2k-1)(k-1)$. Since $q\ge k-1$ we have $m\ge 2k^2-2k+1$, so $2m>2k^2-k$ and
\begin{equation*}
km+k(2k-1)<(k+2)m\le(q+3)m=D+m\,.
\end{equation*}
The quadratic relations have weight $2m+2k<4m\le D$. If $q\ge k$, then $D\ge(k+2)m$ exceeds every weight above. If $q=k-1$, then $m=2k^2-2k+1$ and $D=(k+1)m$, and a relation $f$ of length $k$ has weight at least $D$ exactly when $\delta_f\ge m$. Among the offsets, $\delta(P_4)=2k^2-k=m+k-1$ and $\delta(G_{k-2})=k+(2k-1)(k-1)=m$, while $\delta(G_i)\le m-(2k-1)$ for $i\le k-3$, $\delta(F_i)\le 2k^2-3k+1<m$ and $\delta(P_3)=k<m$.
\end{proof}

So for $q\ge k$ no element of $H$ shares a weight with any~$N_r$, but at the boundary $q=k-1$ the relation $G_{k-2}=be^{k-1}-cd^{k-1}$ has the weight of~$N_0$ and $P_4$ has the weight of~$N_{k-1}$. The map $\psi$ separates them.

\begin{lemma}\label{lem:psi}
$\psi(f)=0$ for every $f\in H$, and $\psi(N_r)=s^{q+1}t^{r}(t^{m}-s)\ne 0$ for $0\le r\le 2k-1$, over every field~$K$.
\end{lemma}

\begin{proof}
For a monomial $M$, $\psi(M)=s^{\len{M}}t^{\delta(M)}$. The two monomials of an element of $H$ have the same length and the same weight by Lemma~\ref{lem:relations}, hence the same offset, so their images cancel. By the same lemma $\psi(U_r)=s^{q+1}t^{m+r}$ and $\psi(V_r)=s^{q+2}t^{r}$. The difference $s^{q+1}t^r(t^m-s)$ has two distinct monomials with coefficients $\pm1$, so it is nonzero in every characteristic.
\end{proof}

\subsection{The relations of~$H$ among themselves}

All elements of $H$ are homogeneous of length $2$ or~$k$, so a redundancy among them can be projected to a single length, where the multipliers of relations of the same length become scalars. Modulo the monomial ideal $(a,c^2,bd)$ the two quadratic relations disappear and the $2k$ relations of length $k$ have disjoint supports.

\begin{lemma}\label{lem:Hnonred}
No element of $H$ lies in the ideal generated by the other elements of~$H$.
\end{lemma}

\begin{proof}
Suppose $f=\sum_{g\in H\setminus\{f\}}c_g\,g$ with $f\in H$ and $c_g\in R$. Taking the homogeneous component of length $\len{f}$ on both sides, we may replace each $c_g$ by its component of length $\len{f}-\len{g}$; this is zero when $\len{g}>\len{f}$ and a scalar when $\len{g}=\len{f}$.

If $f$ is one of the quadratics $P_1,P_2$, only the other quadratic can occur, with a scalar coefficient, because the remaining relations have length $k\ge 3$. But $c^2-bd$ and $bd-ae$ are not proportional.

If $f$ has length~$k$, we obtain
\begin{equation}\label{eq:Hredundancy}
f=c_1P_1+c_2P_2+\sum_{g}\lambda_g\,g\,,
\end{equation}
where $g$ runs over the relations of length $k$ other than~$f$ and $\lambda_g\in K$. Let $Q=(a,c^2,bd)$. Since $Q$ is a monomial ideal, $R/Q$ has as $K$-basis the classes of the monomials outside~$Q$, and the class of a polynomial is determined by its terms with such monomials. Reducing~\eqref{eq:Hredundancy} modulo~$Q$ kills $P_1$, $P_2$, and every monomial containing~$a$. The surviving parts of the relations of length~$k$ are
\begin{equation*}
P_3\mapsto b^k,\qquad P_4\mapsto d^k-ce^{k-1},\qquad F_i\mapsto b^{k-1-i}ce^{i}\ \ (0\le i\le k-2)\,,
\end{equation*}
and $G_i\mapsto b^{k-1-i}e^{i+1}$ for $0\le i\le k-3$, while $G_{k-2}\mapsto be^{k-1}-cd^{k-1}$. None of the monomials shown lies in~$Q$, and they are pairwise distinct: $b^k$ is the only one with $b$-exponent~$k$; the images of $F_i$ are the only ones divisible by $bc$; the images of $G_i$, $i\le k-3$, together with $be^{k-1}$, are the monomials $b^{x}e^{y}$ with $x\ge 1$, $x+y=k$ and $x\le k-1$; and $d^k$, $ce^{k-1}$, $cd^{k-1}$ are distinct and differ from all the others. Hence the images of the $2k$ relations of length $k$ are nonzero with pairwise disjoint supports, so they are linearly independent in $R/Q$, contradicting~\eqref{eq:Hredundancy}.
\end{proof}

\subsection{The full list}

\begin{proposition}\label{prop:nonred}
No element of $H\cup N$ lies in the ideal generated by the other $4k+1$ elements. In particular the $4k+2$ elements of $H\cup N$ are pairwise distinct.
\end{proposition}

\begin{proof}
The second assertion follows from the first, since two equal elements would each lie in the ideal generated by the other. Suppose $f=\sum_{g}c_g\,g$, the sum running over $g\in(H\cup N)\setminus\{f\}$. All the relations are homogeneous for~$\nu$ (Lemma~\ref{lem:relations}), so taking the component of weight $\nu(f)$ we may assume that each $c_g$ is homogeneous of weight $\nu(f)-\nu(g)$, and zero if this is negative. A nonzero homogeneous polynomial has weight $0$, when it is a scalar, or weight at least~$m$. In both cases below, every coefficient $c_{N_j}$ of a relation of~$N$ turns out to have weight less than~$m$, so it is a scalar, and it vanishes unless its weight is~$0$.

\emph{The case $f=N_r$.} For $j\ne r$ the coefficient $c_{N_j}$ has weight $r-j$, which is nonzero and of absolute value at most $2k-1<m$, so $c_{N_j}=0$. Thus $N_r$ lies in the ideal generated by~$H$. Applying $\psi$ and Lemma~\ref{lem:psi} gives $\psi(N_r)=0$, a contradiction.

\emph{The case $f\in H$.} By Lemma~\ref{lem:Hweights}, $\nu(f)<D+m$, so the coefficient of $N_j$ has weight $\nu(f)-D-j<m$; it is therefore a scalar $\lambda_j$, nonzero only if $\nu(f)=D+j$. Applying $\psi$ and Lemma~\ref{lem:psi} gives $0=\sum_j\lambda_j\psi(N_j)$. At most one $\lambda_j$ is nonzero, since the weights $D+j$ are distinct, and $\psi(N_j)\neq 0$; so every $\lambda_j$ vanishes. Then $f$ lies in the ideal generated by $H\setminus\{f\}$, which Lemma~\ref{lem:Hnonred} excludes.
\end{proof}

Nonredundancy is weaker than indispensability, and the stronger property fails. A binomial of $I$ is \emph{indispensable} if, up to sign, it belongs to every system of binomial generators of~$I$~\cite{CKT07}.

\begin{remark}\label{rem:N0}
The relation $N_0$ is not indispensable. Indeed
\begin{equation*}
N_0+e^{q-k+1}\,G_{k-2}=cd^{k-1}e^{q+1-k}-a^{q+2}+be^{q}-cd^{k-1}e^{q-k+1}=be^{q}-a^{q+2}\,,
\end{equation*}
so replacing $N_0$ by $be^{q}-a^{q+2}$ in $H\cup N$ gives another set of $4k+2$ binomials generating~$I$, which is of least cardinality by Theorem~\ref{thm:F} and does not contain~$\pm N_0$.
\end{remark}

\section{From nonredundancy to the minimal number of generators}\label{sec:mu}

Nonredundancy of one generating set does not by itself bound the size of other generating sets: in a general ring an ideal can have nonredundant generating sets of different sizes, as $(x)=(x^2,x-x^2)$ in~$K[x]$ shows. For homogeneous generators of a graded ideal the vector space $I/\mm I$ removes this ambiguity. This is the standard graded Nakayama argument; we give it in the form needed, in every characteristic and for arbitrary, not necessarily binomial or homogeneous, generating sets.

\begin{lemma}\label{lem:graded}
Let $R=K[x_1,\dots,x_n]$ be graded by assigning positive integer weights to the variables, let $\mm=(x_1,\dots,x_n)$, and let $I$ be an ideal generated by homogeneous elements $f_1,\dots,f_N$ such that no $f_j$ lies in the ideal generated by the $f_i$ with $i\ne j$. Then $\dim_K I/\mm I=N$, and every set of generators of $I$ has at least $N$ elements. In particular the least cardinality of a generating set of $I$ is~$N$.
\end{lemma}

\begin{proof}
The quotient $I/\mm I$ is a module over $R/\mm=K$. Its elements are classes of $\sum_ig_if_i$ with $g_i\in R$, and $\sum_ig_if_i\equiv\sum_ig_i(0)f_i$ modulo~$\mm I$, because $g_i-g_i(0)\in\mm$. Hence the classes $[f_1],\dots,[f_N]$ span $I/\mm I$.

They are linearly independent. Suppose $\sum_i\lambda_if_i\in\mm I$ with $\lambda_j\ne 0$, and let $\sigma$ be the weight of~$f_j$. Every element of $\mm I$ has the form $\sum_ih_if_i$ with $h_i\in\mm$, because a product $xy$ with $x\in\mm$ and $y=\sum_ig_if_i\in I$ equals $\sum_i(xg_i)f_i$. Since the variables have positive weights, the elements of $\mm$ are exactly the polynomials with zero component of weight~$0$. Taking components of weight $\sigma$ in $\sum_i\lambda_if_i=\sum_ih_if_i$ gives
\begin{equation*}
\sum_{i:\,\nu(f_i)=\sigma}\lambda_if_i=\sum_i\bigl(h_i\bigr)_{\sigma-\nu(f_i)}\,f_i\,,
\end{equation*}
where $(h)_\tau$ denotes the component of weight $\tau$ and $\nu(f_i)$ the weight of~$f_i$. On the right only the $i$ with $\nu(f_i)<\sigma$ contribute, since $(h_i)_\tau=0$ for $\tau\le 0$. Solving for $\lambda_jf_j$ expresses $f_j$ as a combination of the $f_i$ with $i\ne j$, contrary to the hypothesis. Hence $\dim_KI/\mm I=N$.

If $G$ is any set of generators of $I$, the classes of its elements span $I/\mm I$ by the argument of the first paragraph, so $G$ has at least $N$ elements. The $f_i$ are pairwise distinct, since $f_i=f_j$ with $i\ne j$ would put $f_j$ in the ideal generated by~$f_i$, so $\{f_1,\dots,f_N\}$ is a generating set with exactly $N$ elements.
\end{proof}

\begin{proof}[Proof of Theorem~\ref{thm:F}]
Grade $R=K[a,b,c,d,e]$ by the weights $\nu$, which are positive. The $4k+2$ elements of $H\cup N$ are homogeneous by Lemma~\ref{lem:relations}, generate $I$ by Proposition~\ref{prop:generation}, and form a nonredundant list by Proposition~\ref{prop:nonred}. Lemma~\ref{lem:graded} gives $\mu(I)=4k+2=2w+2$.
\end{proof}

\section{Formal verification}\label{sec:lean}

Theorem~\ref{thm:F} has been verified in Lean~4 with Mathlib.\footnote{Toolchain \texttt{leanprover/lean4:v4.35.0-rc3}; Mathlib commit \texttt{c55e6e786f49471c72fbddbec5415808896aec1e}.} The formal statement, named \texttt{FiveGenBetti.family\_spanRank\_eq}, asserts for every field $K$ and all natural numbers $k\ge 3$ and $q\ge k-1$ that the kernel of the map \texttt{MvPolynomial.aeval} sending the five variables to the powers of $t$ with exponents $m,m+1,m+k,m+2k-1,m+2k$ has \texttt{Submodule.spanRank} equal to $4k+2$. Since \texttt{spanRank} is the least cardinality of a generating set, this is exactly the statement $\mu(I)=4k+2$. The formal proof follows the route of Sections~\ref{sec:family}--\ref{sec:mu}, except that it proves Lemma~\ref{lem:Hnonred} by a divisibility argument with witness monomials where the paper projects modulo $(a,c^2,bd)$. It uses only the axioms \texttt{propext}, \texttt{Classical.choice} and \texttt{Quot.sound}, has been replayed by the Lean kernel, and passed the Palomar registry's comparator script when run in continuous integration at commit \texttt{0223cfc} (statement comparison, axiom check, and replay by the independent kernels con-ron and nanoda). Conjecture~\ref{conj:W} is not proved, formally or otherwise. The source is in the repository \texttt{github.com/benkeene/five-generator-betti}, available from the author \red{[visibility]}. The proofs were developed with substantial assistance from GPT-6 Astra (OpenAI) and Claude models (Anthropic), under an adversarial review protocol; the formal verification, not that review, is the guarantee of correctness.

\enlargethispage{3\baselineskip}
\bibliographystyle{plain}
\bibliography{refs}

\end{document}
